\beginsong{Des signes de fraternité}[by={Laurent Grzybowski}]
\textnote{Capo 2}
\beginverse
Ça \[Em]chauffe un peu partout sur \[C]la \[D]plan\[Em]ète,
Au \[Em]c{\oe}ur des mégapoles, on se fl\[C]ingue à \[D]tout \[Em]va.
La \[Em]raison du plus fort est la \[C]plus\[D] bê\[Em]te,
\[Em]Mais par-dessus tout ça, \[C]regarde bien tu ver\[D]ras :
\endverse
\beginchorus
Des \[C]signes de frat\[D]ern\[G]ité,
\[Em]Dans tous les\[C] recoins\[G] de \[D]la terre.
\[C]Des signes de fra\[D]tern\[Em]ité,
\[Am]Sur les barbelés de \[B7]la guerre.
\[C]Des signes de \[D]frate\[G]rnité.
\[Em]Des \[C]signes de fra\[D]ternit\[Em]é.
\endchorus
\beginverse
Tou\[Em]jours chacun pour soi, cha\[C]cun \[D]sa \[Em]route,
On \[Em]fait des hécatombes au vol\[C]ant de \[D]la \[Em]vie.
Pour \[Em]passer le premier quoi \[C]qu'il \[D]en \[Em]coûte,
\[Em]Mais par-dessus tout ça, \[C]regarde bien tu ver\[D]ras :
\endverse
\beginverse
Part\[Em]out les vieux démons rev\[C]iennent \[D]en \[Em]force,
Et \[Em]braillent des slogans pour ban\[C]nir l'ét\[D]ran\[Em]ger.
D'un \[Em]quartier d'un pays, au \[C]c{\oe}ur \[D]des \[Em]hommes,
\[Em]Mais par-dessus tout ça, \[C]regarde bien tu ver\[D]ras :
\endverse
\beginverse
Des \[Em]peuples cousus \[C]dans la \[D]mi\[Em]sère
Et \[Em]croisent l'abondance aux écr\[C]ans des \[D]nan\[Em]tis.
De \[Em]quoi désespérer \[C]la vie \[D]en\[Em]tière,
\[Em]Mais par-dessus tout ça, \[C]regarde bien tu ver\[D]ras :
\endverse
\beginverse
\[Em]Devant la religion du terro\[C]risme, \[D] \[Em]
\[Em]Barbus de Barbarie, fous d'un dieu-Atti\[C]la. \[D] \[Em]
\[Em]Le monde a des relents d'appoca\[C]lypse, \[D] \[Em]
\[Em]Mais par-dessus tout ça, \[C]regarde bien tu ver\[D]ras :
\endverse
\beginverse
\[Em]Regarde tu verras, le monde \[C]change, \[D] \[Em]
\[Em]Il y a des mains tendues et des c{\oe}urs grands ou\[C]verts \[D] \[Em]
\[Em]Qui font se mélanger les diffé\[C]rences, \[D] \[Em]
\[Em]Dans les rues de ta vie, tu vois fleurir aujourd'\[C]hui ! \[D] \[Em]
\endverse
\endsong
